\section*{Bloque B. Aplicaciones cuantitativas y optimización}

\noindent\textit{Preguntas 31 a 65. Requieren calcular el ingreso marginal del producto, el costo marginal del factor, equilibrios competitivos y de monopsonio, efectos de salarios mínimos, elasticidades y resultados de negociación.}

\noindent\textit{Para las preguntas 31 a 45, considere un mercado de trabajo donde la demanda de trabajo está dada por $w=70-2L$ y la oferta de trabajo por $w=10+L$. El salario se expresa en unidades monetarias y $L$ en unidades de trabajo.}

\q{La cantidad de trabajo del equilibrio competitivo es:}
\begin{opts}
  \item 20
  \item 15
  \item 25
  \item 30
\end{opts}

\q{El salario del equilibrio competitivo es:}
\begin{opts}
  \item 20
  \item 25
  \item 35
  \item 30
\end{opts}

\q{Si una empresa actúa como monopsonista, la cantidad de trabajo que maximiza su beneficio es:}
\begin{opts}
  \item 20
  \item 15
  \item 25
  \item 10
\end{opts}

\q{El salario que paga el monopsonista es:}
\begin{opts}
  \item 20
  \item 30
  \item 15
  \item 25
\end{opts}

\q{El costo marginal del factor trabajo en el óptimo del monopsonista es:}
\begin{opts}
  \item 25
  \item 40
  \item 30
  \item 55
\end{opts}

\q{El ingreso marginal del producto del trabajo en el óptimo del monopsonista es:}
\begin{opts}
  \item 40
  \item 25
  \item 30
  \item 55
\end{opts}

\q{La explotación monopsonística, medida como la diferencia entre el MRP y el salario, es:}
\begin{opts}
  \item 5
  \item 10
  \item 20
  \item 15
\end{opts}

\q{La diferencia de empleo entre el equilibrio competitivo y el monopsonio es:}
\begin{opts}
  \item 2
  \item 5
  \item 10
  \item 15
\end{opts}

\q{La diferencia salarial entre el equilibrio competitivo y el monopsonio es:}
\begin{opts}
  \item 2
  \item 10
  \item 15
  \item 5
\end{opts}

\q{Si se fija un salario mínimo igual al salario competitivo, el empleo del monopsonista será:}
\begin{opts}
  \item 20
  \item 15
  \item 25
  \item 30
\end{opts}

\q{El salario mínimo que lleva al monopsonista a contratar la cantidad competitiva es:}
\begin{opts}
  \item 20
  \item 25
  \item 35
  \item 30
\end{opts}

\q{El costo marginal del factor trabajo cuando $L=20$ es:}
\begin{opts}
  \item 30
  \item 40
  \item 50
  \item 60
\end{opts}

\q{La renta económica de los trabajadores en el equilibrio competitivo es:}
\begin{opts}
  \item 100
  \item 200
  \item 300
  \item 400
\end{opts}

\q{La renta económica de los trabajadores en el monopsonio es:}
\begin{opts}
  \item 100
  \item 150
  \item 112.5
  \item 225
\end{opts}

\q{La pérdida de eficiencia del monopsonio respecto del equilibrio competitivo es:}
\begin{opts}
  \item 12.5
  \item 25
  \item 50
  \item 100
\end{opts}

\q{Si la oferta de trabajo es $w=10+L$, el costo total del trabajo es:}
\begin{opts}
  \item $10L$
  \item $10L+L^2$
  \item $10+L$
  \item $20L$
\end{opts}

\q{Con la oferta de trabajo $w=10+L$, el costo marginal del factor trabajo es:}
\begin{opts}
  \item $10$
  \item $10+L$
  \item $20+2L$
  \item $10+2L$
\end{opts}

\q{Si el producto marginal del trabajo es 10 y el precio del producto es 5, el ingreso marginal del producto del trabajo es:}
\begin{opts}
  \item 10
  \item 15
  \item 50
  \item 5
\end{opts}

\q{Si el producto marginal del trabajo es 8 y el precio del producto es 10, el ingreso marginal del producto del trabajo es:}
\begin{opts}
  \item 18
  \item 10
  \item 80
  \item 8
\end{opts}

\q{Si el MRP es 60 y el salario es 40, la explotación monopsonística es:}
\begin{opts}
  \item 10
  \item 40
  \item 60
  \item 20
\end{opts}

\q{Si la elasticidad de la demanda de trabajo es $-2$ y el salario aumenta 10\%, el empleo cae aproximadamente:}
\begin{opts}
  \item 20\%
  \item 10\%
  \item 5\%
  \item 2\%
\end{opts}

\q{Si la elasticidad de la oferta de trabajo es 0.5 y el salario aumenta 10\%, la oferta de trabajo aumenta:}
\begin{opts}
  \item 10\%
  \item 0.5\%
  \item 5\%
  \item 20\%
\end{opts}

\q{Si el salario sube de 20 a 22 y la cantidad de trabajo de 100 a 110, la elasticidad arco de la oferta de trabajo es aproximadamente:}
\begin{opts}
  \item 0.5
  \item 1.0
  \item 2.0
  \item 0.1
\end{opts}

\q{Si la demanda de trabajo es $L=100-2w$, la elasticidad de la demanda en $w=20$ es aproximadamente:}
\begin{opts}
  \item $-1.5$
  \item $-2.0$
  \item $-0.67$
  \item $0.5$
\end{opts}

\q{Si la demanda de trabajo es $L=100-2w$ y la oferta es $L=2w-20$, el salario de equilibrio es:}
\begin{opts}
  \item 30
  \item 20
  \item 40
  \item 50
\end{opts}

\q{En el equilibrio anterior, la cantidad de trabajo es:}
\begin{opts}
  \item 20
  \item 40
  \item 60
  \item 80
\end{opts}

\q{Si el gobierno fija un salario mínimo de 35 en ese mercado competitivo, el empleo será:}
\begin{opts}
  \item 40
  \item 50
  \item 30
  \item 20
\end{opts}

\q{El desempleo generado por el salario mínimo de 35 es:}
\begin{opts}
  \item 10
  \item 30
  \item 40
  \item 20
\end{opts}

\q{Si el salario mínimo se fija en 20, por debajo del equilibrio, el empleo efectivo es:}
\begin{opts}
  \item 20
  \item 30
  \item 60
  \item 40
\end{opts}

\q{En un monopsonio, si el salario mínimo se fija en el salario competitivo, el empleo:}
\begin{opts}
  \item aumenta hasta el nivel competitivo.
  \item disminuye.
  \item no cambia.
  \item se vuelve cero.
\end{opts}

\q{Si el sindicato fija un salario por encima del competitivo en un mercado competitivo, el empleo:}
\begin{opts}
  \item aumenta.
  \item no cambia.
  \item disminuye.
  \item se vuelve cero.
\end{opts}

\q{En un monopolio bilateral, el resultado del salario depende de:}
\begin{opts}
  \item solo la demanda de trabajo.
  \item solo la oferta de trabajo.
  \item el poder de negociación de las partes.
  \item únicamente el gobierno.
\end{opts}

\q{En una negociación entre un sindicato y una empresa, el salario resultante se ubicará:}
\begin{opts}
  \item dentro de un rango delimitado por el salario de reserva y el salario máximo que la empresa está dispuesta a pagar.
  \item en cero.
  \item por encima del MRP.
  \item por debajo de la oferta de trabajo.
\end{opts}

\q{La curva de contrato en la negociación salarial se obtiene cuando:}
\begin{opts}
  \item la demanda iguala a la oferta.
  \item las curvas de indiferencia del sindicato y las curvas de isobeneficio de la empresa son tangentes.
  \item el MRP es cero.
  \item el salario es cero.
\end{opts}

\q{Si el sindicato maximiza la masa salarial $wL$ y la demanda de trabajo es $L=100-2w$, el salario que maximiza $wL$ es:}
\begin{opts}
  \item 20
  \item 30
  \item 25
  \item 40
\end{opts}
